% problem-000075 9 有机合成与机理
\section*{9. 有机合成与机理}
\textit{下列为分问。}

\subsection*{9-1}
画出G的结构简式。

\subsection*{9-2}
研究发现，当使⽤超酸三氟甲磺酸代替浓硫酸后，反应的产物如下图左所⽰。研究还表明，化合物G在三氟甲磺酸作⽤下也可以转化为该化合物。画出化合物G转化为此产物过程中所形成的中间体。提⽰：同⼀物种只须画出其中⼀个主要共振式。

（图：）

\subsection*{9-3}
参照以上实验结果，如果希望制备以下内酯，画出所⽤原料H的结构简式（说明：LA即Lewis酸）：

（图：）

\subsection*{9-4}
室温下五元环状化合物克酮酸(croconic acid, $\mathrm { C 5 H _ { 2 } O _ { 5 } ) }$ 在浓硫酸的作⽤下先与等量的苯反应形成中间体I$\mathrm { ( C _ { 1 1 } H _ { 6 } O _ { 4 } ) }$ ；I继续在苯中反应⽣成J $\left( \mathrm { C } _ { 2 3 } \mathrm { H } _ { 1 6 } \mathrm { O } _ { 3 } \right)$ 。画出I和J的结构简式。
